\section{Related Work}
\label{sec:related}

\subsection{Automatic Fusion of Dependent Reductions: Neptune and RedFuser}
\label{sec:related-prior}

近期有两个编译器能够自动融合依赖归约，它们是与本文最接近的工作。
\emph{Neptune}~\citep{neptune} 先按不存在依赖的方式融合归约，再通过符号求解得到修复函数 $h(t,r,r')$，将使用较早归约的过时值 $r$ 计算出的部分结果 $t$，转换为对应新值 $r'$ 的结果。
融合通过用户在调度模板中指定的滚动更新和 split-K 调度实现，并由 Triton 编译~\citep{triton}；根据 Neptune 的结论，其后续工作 Nautilus~\citep{nautilus} 将调度自动化。
\emph{RedFuser}~\citep{redfuser} 检查每个归约体能否分解为 $G(x)\otimes H(d)$，即数据的函数与较早结果的函数的组合，其中 $\otimes$ 由归约类型决定；随后推导通过 $H$ 重新缩放的融合形式与增量形式，并生成 TileLang 内核~\citep{tilelang}。
两者所述分析都根据已有归约结果和变化的生产者值推导修正，不综合新的聚合定义。
两个系统也都评估 GPU 性能；\sys 目前尚未进行此类性能评估。

\textbf{Their conditions.}
我们比较两者对消费者归约的代数条件：归约运算为 $\oplus\in\{\Sigma,\max\}$（$\min$ 为其对偶），归约体为 $g(r,c)$，上下文为 $r$，不依赖上下文的值为 $c$。
\begin{itemize}
\item[$\Cls{N}$:] Neptune 的条件~\citep[Eq.~3--4]{neptune}：存在某个 $h$，使得对所有 $r,r',c,t,t'$，有 $h(g(r,c),r,r')=g(r',c)$ 和 $h(t\oplus t',r,r')=h(t,r,r')\oplus h(t',r,r')$。
\item[$\Cls{R}$:] RedFuser 的分解条件~\citep[Eq.~4, Tab.~1]{redfuser}：当 $\oplus=\Sigma$ 时，$g(r,c)=G(c)\cdot H(r)$；当 $\oplus=\max$ 时，$g(r,c)=G(c)+H(r)$。
\item[$\Cls{H}$:] \cref{lem:regroup} 的条件，其中包含有限个不依赖上下文的 $G_k$：
\[
g(r,c)=\sum_{k=1}^{K}a_k(r)\,G_k(c)\ \text{ when }\oplus=\Sigma,\qquad
g(r,c)=\max_{k=1}^{K}\psi_k(r,G_k(c))\ \text{ when }\oplus=\max,
\]
其中每个 $\psi_k(r,\cdot)$ 都单调不减。
\end{itemize}
RedFuser 的合并推导还要求旧上下文因子在其修正运算下可逆~\citep[\S3.2.2]{redfuser}。
下文的类别 $\Cls{R}$ 描述归约体的分解，这一执行条件另行讨论。

\begin{theorem}[Containment]
\label{thm:containment}
(a) $\Cls{R}\subseteq\Cls{H}$。
(b) $\Cls{N}$ 中的每个归约都属于 $\Cls{H}$，但在 $\oplus=\Sigma$ 时，要求其修复函数 $h(\cdot,r,r')$ 至少在某一点连续。
(c) 两个包含关系都是严格的：$\sum_j (c_j-r)^2$ 和 $\max_j |c_j-r|$ 属于 $\Cls{H}$，但既不属于 $\Cls{N}$，也不属于 $\Cls{R}$。
\end{theorem}

\begin{proof}[Proof sketch]
(a) $G(c)H(r)$ 对应 $K=1$、$a_1=H$ 的情形；$u\mapsto u+H(r)$ 非递减。
(b) 对于 $\oplus=\Sigma$，第二个等式使 $h(\cdot,r,r')$ 可加，因此一旦它在某一点连续，就有线性形式 $h(t,r,r')=\lambda(r,r')\,t$（Cauchy 方程~\citep{kuczma}）。
第一个等式将 $r$ 固定为某个 $r_0$，得到单项形式 $g(r',c)=\lambda(r_0,r')\,g(r_0,c)$。
对于 $\oplus=\max$，第二个等式使 $h(\cdot,r,r')$ 非递减，第一个等式给出 $g(r',c)=h(g(r_0,c),r_0,r')$，它是 $G(c)=g(r_0,c)$ 的非递减函数。
(c) $(c-r)^2=c^2-2rc+r^2$ 和 $|c-r|=\max(c-r,\,r-c)$ 均属于 $\Cls{H}$。
改变上下文会区分被这些归约体映射为相同部分值的输入，因此仅凭该值无法确定 $\Cls{N}$ 所需的修正。
这两种归约体也不具有 $\Cls{R}$ 要求的单因子或加法分离形式。
补充材料附录 A 给出完整证明。
\end{proof}

定理所比较的条件将消费者的上下文与输入值视为独立参数。
被这些条件排除的归约体需要消费者部分结果之外的信息；但所在的完整程序可能已经保留了部分此类信息，或能建立结果之间的额外关系。
Neptune 报告其求解器无法推导 Welford 方差，是因为未利用均值、方差和新输入样本之间的关系~\citep[\S4.2]{neptune}。
RedFuser 报告了方差和转动惯量基准~\citep[App.~A.6]{redfuser}，但未解释其中心化平方归约体如何与所述分解检查相容。
因此，包含关系针对陈述的代数条件，而下文的程序比较针对论文报告的工作负载。
识别能力还取决于符号过程：Neptune 和 RedFuser 使用 SymPy，\sys 使用 \cref{tab:rules} 中的重写规则。
连续性假设排除了病态的不可测可加修复函数。

\textbf{Their programs.}
为比较实现，我们将两篇论文融合的 \nPriorPrograms 个程序改写为 TE，并交给快照 \code{6c46fb1b02} 中的 \sys 处理（补充材料附录 F）。
该版本成功融合其中 \nPriorFused 个，包括将掩码或偏置作为张量提供后的所有注意力变体、RMSNorm 后接最大值、Performer 的三种归约模式、MLA 解码、方差和转动惯量。
六个程序使用了 \cref{fig:grammar} 语言之外的结构：用 \code{if\_then\_else} 构造掩码、由循环索引计算 ALiBi、使用 $\tanh$ 的软截断、通过 \code{n // G} 索引的分组查询，以及同时在两个轴上的归约。
缩放 L2 范数和 FP8 量化后接 GEMM 得到了经验证的归约器，但在极端输入族上未通过当时版本的数值检查；它们在更新后的检验方法下的结果仍待测量。
另有两个程序未通过状态选择：嵌套求和示例存在未解决的定义域义务，下文附录 A.1 程序所需的辅助状态则超出元组大小限制。
Neptune 在 L2 范数中保留缩放，并除以运行中的最大值；在全零前缀之后，该值为零，而论文未讨论这种情况。

\textbf{Contexts without an inverse.}
RedFuser 的推导使用 $H(d_{\mathrm{old}})^{-1}\otimes H(d_{\mathrm{new}})$ 重新缩放部分结果。
附录 A.1 将 $H$ 的不可逆取值替换为单位元，得到 $H'$，并在新旧两个因子中均使用 $H'$~\citep[Eq.~27--28]{redfuser}。

\label{ex:a1}
当局部的零因子已经抹去某项贡献时，这种替换无法从局部结果中恢复该贡献。
有效的合并必须保留足以重构被抹去贡献的信息，例如增加一个辅助状态。
在 \sys 的保留策略下，该辅助状态必须与有定义的原始结果共同保留；如果所得元组超出大小限制，该链保持未融合。
补充材料附录 D 给出针对这一修正后合并公式的数学反例。

\Cref{tab:capabilities} 汇总了比较结果，也包括前人系统领先的方面。

\begin{table}[t]
\caption{Comparison of stated analyses and reported capabilities. Algebraic rows refer to the stated body conditions; workload rows refer to reported programs. \cmark: supported. \xmark: excluded by that condition, or not done. \pmark: partly. \nr: not reported. Page numbers refer to the papers.}
\label{tab:capabilities}
\footnotesize
\setlength{\tabcolsep}{3pt}
\renewcommand{\arraystretch}{1.1}
\begin{tabular}{@{}>{\raggedright\arraybackslash}p{0.4\linewidth}>{\raggedright\arraybackslash}p{0.19\linewidth}>{\raggedright\arraybackslash}p{0.19\linewidth}>{\raggedright\arraybackslash}p{0.17\linewidth}@{}}
\toprule
& \textbf{Neptune} & \textbf{RedFuser} & \textbf{\sys}\\
\midrule
Sum with one multiplicative context (softmax, $P\cdot V$) & \cmark{} Eq.~3--5, p.~7 & \cmark{} Eq.~4, Tab.~1 & \cmark{}\\
Max with an additive or a positive-scale context & \cmark{} p.~13 & \pmark{} additive only, Tab.~1 & \cmark{}\\
Bodies $(c-r)^2$ and $|c-r|$ under the stated conditions & \xmark{} p.~10--11; \cref{thm:containment} & \xmark{} \cref{thm:containment} & \pmark{} subject to tuple size limit\\
Synthesis of new aggregate definitions & \xmark{} & \xmark{} & \cmark{}\\
Join proven equal to the program over $\R$ & \cmark{} Thm.~4.1, 4.4 & \pmark{} derivation & \cmark{} \cref{thm:correct}\\
Noninvertible multiplicative context factor & \nr & \xmark{} App.~D & \pmark{} verified guards\\
Accuracy relative to the original decides fusion & \xmark{} measured after, App.~F & \nr & \cmark{} \cref{def:gate}\\
Masked, ALiBi, soft-cap, GQA attention as written & \cmark{} Tab.~1, p.~16 & \pmark{} MHA, MLA & \xmark{} frontend\\
Reducers other than sum, max, min & \cmark{} associative formulation & \cmark{} top-k; transformed product & \xmark{}\\
Scheduling, GPU code generation, performance & \cmark{} Triton & \cmark{} TileLang & \xmark{} \cref{sec:gpu}\\
\bottomrule
\end{tabular}
% src: Neptune PDF pages 7, 10-11, 13, 16, App. F; RedFuser PDF Eq. 4, Tab. 1 p.7, A.1 p.15, A.6 p.17;
%      Hapax: data/prior.json, data/fused_ops.json; theorems of this paper.
\end{table}

\subsection{Other Related Work}

\textbf{Hand-derived single passes.}
Welford 更新~\citep{welford}、其加权形式~\citep{west}、\citet{chan} 研究的成对合并，以及 \citet{pebay} 的任意阶矩合并，都是面向统计计算的单遍算法。
在线 softmax~\citep{milakov}、FlashAttention~\citep{flashattention,flashattention2}、cut cross-entropy~\citep{cce} 和 FlashSinkhorn~\citep{flashsinkhorn} 是面向机器学习的单遍算法；缩放欧几里得范数~\citep{blue} 则通过维护缩放来控制中间值的范围。
这些算法都针对特定算子手工推导；\sys 从算子定义推导 softmax、注意力和方差归约器。

\textbf{Tensor compilers and superoptimisers.}
Flashlight~\citep{flashlight} 利用指数函数的同态性质，将注意力变体重写为在线形式。
Halide 的 \code{rfactor}~\citep{rfactor} 和 TVM~\citep{tvm,tensorir} 可以并行化满足结合律的归约，包括元组归约器，但不能将一个归约与读取其最终值的后续归约融合。
Mirage~\citep{mirage} 搜索多层张量算子图，并在包含多线性算子、除法和每条路径至多一次指数运算的片段上，对候选进行概率验证，再用浮点测试过滤数值不稳定的候选；最大值不在该片段中。
STENSO~\citep{stenso} 和 Prism~\citep{prism} 通过符号方法搜索等价张量程序。
\sys 将综合限制为在固定源归约结果的基础上寻找辅助状态。
上下文修正表达式提供文法终结项，完整参考恢复提供综合规约。
这将表达式搜索与合并操作的构造和验证分离。

\textbf{Program calculation and synthesis.}
列表同态~\citep{bird87} 和第三同态定理~\citep{gibbons96} 刻画了列表函数何时具有分治形式；\citet{morita07} 从两个顺序程序构造这种形式，而 catamorphism 等递归模式~\citep{meijer91} 给出了此类推导背后的演算规律。
元组化消除对同一数据结构的多次遍历~\citep{bird84,pettorossi87,hu96,hu97}。
Farzan 和 Nicolet 通过展开循环体并收集辅助累加量来综合分治程序~\citep{farzan17,farzan19,farzan21}。
AutoLifter 将提升问题分解为辅助值和组合器的较小综合任务~\citep{autolifter}；SuFu 则综合压缩函数和标量计算，以消除中间数据结构~\citep{sufu}。
\citet{morihata21} 通过半环多项式并行化归约。
Opera~\citep{opera} 从离线算法推导每次处理一个元素并更新状态的在线算法，在 51 个基准中解决了 50 个，唯一失败的是峰度。
% src: notes/evidence/opera.txt lines 913, 934 (50 of 51; kurtosis)
这些工作为顺序循环和递归程序发展了综合与演算方法。
\sys 聚焦于单个张量轴上的归约，因此能够通过有限展开建立充分条件，并推导多重集上的可交换合并。
在这一范围内，上下文分离给出合并操作已知的参考表示。
要求恢复这一表示，就能消除独立的组合器综合任务：搜索选择辅助表达式，经过验证的恢复关系则将参考状态的合并转化为维护状态的合并。
固定源状态的定义进一步限制未知表示，上下文修正系数则为其提供候选表达式。
这些限制保留了可用于构造局部中心化或缩放更新的量；生成实现的数值表现另行评估。

\textbf{Aggregation.}
在数据库中，如果部分结果可以直接组合，则聚合是\emph{分配型}的；如果固定元组的分配型聚合足以计算结果，则它是\emph{代数型}的~\citep{datacube}。分布式系统通过树或 gossip 协议计算此类聚合~\citep{jesus}。
不依赖上下文的参考构造使归约链成为代数型聚合；\sys 随后推导所维护的状态元组及其合并操作。
Ink~\citep{ink} 验证或否定给定的用户自定义数据框聚合是否为同态，并在建立该性质时构造合并运算。
其问题是组合给定聚合的结果，而 \sys 发现组合源归约结果所需的辅助状态。
Ink 允许不可交换的合并；\sys 则基于多重集语义推导可交换的合并。

\textbf{Floating-point accuracy.}
经过验证的浮点编译保留数值推理所依赖的源程序语义~\citep{boldo_fp}。
数值分析也表明，中间表示的选择会影响误差传播~\citep{higham_decomposition}；平移后的 log-sum-exp 和 softmax 公式的分析给出了具体例子~\citep{blanchard}。
这些结果为我们围绕源归约定义限制状态选择、并将合并表达为局部上下文修正提供了依据。
Herbie~\citep{herbie} 搜索更准确的表达式形式，FPTaylor~\citep{fptaylor} 则计算严格的舍入误差界。
\sys 为融合发现辅助状态，并使用 MPFR 求值和采样误差容限，以源程序为基准检查所得实现（\cref{def:gate}）。
对生成的合并及其归约树进行舍入误差分析，可以提供超出这些测量的保证。
